\chapter{Podsumowanie}
\label{ch:podsumowanie}

Celem niniejszej pracy było zaprojektowanie i~implementacja kompletnego systemu do przetwarzania i~archiwizacji
strumieniowych danych finansowych z~rynku kryptowalut w~środowisku lokalnym (\textit{on-premises}).
Założenia te zostały zrealizowane w~pełni --- powstał działający system oparty na architekturze mikrousługowej,
który integruje wszystkie zaplanowane komponenty i~spełnia postawione wymagania funkcjonalne oraz niefunkcjonalne.

W~ramach pracy zaimplementowano niezawodny serwer strumieniowy utrzymujący połączenia z~WebSocket giełdy Binance,
automatycznie wznawiający połączenie w~przypadku awarii.
Zbudowano kanał dystrybucyjny oparty na~NATS JetStream, zapewniający buforowanie zdarzeń oraz gwarantowane
dostarczenie komunikatów nawet w~przypadku chwilowej niedostępności usług konsumenckich.
Utworzono procesy ETL dzielące dane na~zbiory surowe (\textit{raw}) i~oczyszczone (\textit{curated}) oraz wzbogacające
je o~metadane analityczne.
Wdrożono mechanizm pamięci warstwowej (\textit{tiered storage}) automatycznie przenoszący dane historyczne
z~operacyjnej bazy MongoDB do~archiwalnego magazynu obiektowego MinIO, optymalizując zasoby warstwy operacyjnej
przy~zachowaniu dostępności danych historycznych.
Zaimplementowano interfejs GraphQL umożliwiający wykonywanie elastycznych zapytań o dane bieżące z~zastosowaniem filtrów
i~projekcji pól.
Zrealizowano bota platformy Telegram realizującego funkcję powiadamiania użytkowników o~zdarzeniach rynkowych
w~czasie rzeczywistym.

System został pomyślnie uruchomiony i~zweryfikowany pod kątem poprawności działania poszczególnych komponentów oraz
przepływu danych między nimi.
Testy manualne opisane w~rozdziale~\ref{ch:testy} potwierdziły prawidłowe funkcjonowanie wszystkich modułów --- od
pobierania danych przez przetwarzanie i~archiwizację, aż po udostępnianie przez API oraz wysyłanie powiadomień.
Weryfikacja mechanizmu migracji danych między warstwami wykazała poprawne przenoszenie rekordów z~MongoDB do MinIO
bez utraty informacji.

Projekt udowodnił, że możliwe jest zbudowanie kompletnego, autonomicznego systemu przetwarzania danych rynkowych
w~środowisku lokalnym z~wykorzystaniem otwartych technologii.
Architektura mikrousługowa sterowana zdarzeniami (\textit{Event-Driven Architecture}) okazała się odpowiednia dla
tego typu zastosowań, zapewniając luźne powiązanie między komponentami, odporność na awarie oraz możliwość
niezależnego skalowania i~rozwoju poszczególnych modułów.

System można rozwijać w~wielu kierunkach, przede wszystkim poprzez centralizację logów z~zastosowaniem
identyfikatorów korelacji oraz wspólnego backendu typu Loki czy~ELK, co~ułatwiłoby diagnozę problemów w~systemach
rozproszonych.
Istotnym uzupełnieniem byłaby ekspozycja metryk technicznych i~domenowych do~Prometheus oraz budowa dashboardów
w~Grafanie, umożliwiających monitorowanie opóźnień, wydajności procesów ETL oraz stanu konsumentów JetStream.
Wdrożenie śledzenia rozproszonego z~wykorzystaniem OpenTelemetry pozwoliłoby na~wizualizację pełnej ścieżki
przetwarzania zdarzeń między mikroserwisami.
W~warstwie aplikacyjnej warto rozważyć rozszerzenie API GraphQL o~dostęp do~danych historycznych z~MinIO poprzez
mechanizm federacji oraz wzbogacenie modułu analitycznego o~bardziej zaawansowane metryki, takie jak~zmienność
w~oknach kroczących, wykrywanie anomalii czy~budowę modeli predykcyjnych w~ramach potoku uczenia maszynowego.
Dodatkowo, wprowadzenie alertów operacyjnych informujących o~problemach działania systemu oraz~skalowanie do~pełnej
listy symboli giełdowych wymagałoby mechanizmów fragmentacji i~lepszego zarządzania backpressure.
Zaprezentowana architektura stanowi solidny fundament do realizacji tych rozszerzeń.

Implementacja systemu pozwoliła na~praktyczną weryfikację wybranych technologii.
Język Go~okazał się niezwykle efektywnym wyborem do~budowy usług sieciowych i~przetwarzania strumieniowego, oferując
niskie zużycie zasobów przy~wysokiej wydajności.
Jasna składnia języka oraz wbudowana obsługa kontekstów i~anulowania operacji znacząco ułatwiły implementację
niezawodnych komponentów, takich jak serwer strumieniowy czy~konsumenci JetStream.
Baza MongoDB sprawdziła się jako warstwa operacyjna, zapewniając prostotę w~modelowaniu dokumentów oraz elastyczność
schematów, co~jest szczególnie istotne w~kontekście danych o~zmiennej strukturze.
W~porównaniu z~relacyjnymi bazami danych, MongoDB wymagała mniejszego nakładu na~konfigurację i~zarządzanie migracjami,
co~przyspieszyło iteracje rozwojowe.
Istnotnym ograniczeniem okazała się sytuacja wokół MinIO --- w~trakcie realizacji projektu producent zaprzestał
rozwijania wersji open~source na~rzecz w~pełni komercyjnego rozwiązania enterprise.
W~konsekwencji wykorzystana wersja nie~otrzymuje już aktualizacji, wsparcia technicznego, a~oficjalna dokumentacja
została usunięta, co~może stanowić wyzwanie w~kontekście długoterminowego utrzymania systemu.
Pomimo tego baza MinIO spełniła swoją rolę w~ramach demonstratora, lecz~w~środowisku produkcyjnym zasadne byłoby rozważenie
alternatyw, takich jak SeaweedFS czy~bezpośrednie wykorzystanie rozwiązań chmurowych kompatybilnych z~API S3.

Podsumowując, wszystkie cele szczegółowe określone we wstępie zostały osiągnięte.
System działa zgodnie z~założeniami, spełnia wymagania funkcjonalne i~niefunkcjonalne, a~przeprowadzone testy
potwierdzają jego poprawność i~stabilność.

